% ============================================================
%  EXPERIMENT 3 — MAIN SECTION  (training & transfer; synthetic + real pastcasting)
%  Place after the Experiment 2 section.
% ============================================================

% ------------------------------------------------------------------
\section{Experiment 3: Training and Transfer Across Synthetic and Real Forecasting Worlds}
\label{sec:exp3}
% ------------------------------------------------------------------

Experiments~1 and~2 treated inductive world-modeling as a \emph{fixed} property of an agent---a
disposition we probe but do not change. They left the third rung of our diagnostic
(our third criterion, \emph{inductive transfer}) untouched, and they left a
sharp open question. In Experiment~2 the strongest agents recovered the hidden gain from a single
observation but \emph{plateaued}: they reached, yet never cleared, a well-posed same-information
forecaster, and were overtaken as evidence accumulated (Table~\ref{tab:exp2-skill}). Experiment~3
asks whether \emph{forecasting training} can move that ceiling---and, crucially, whether any gain is a
\emph{general} improvement in inductive world-modeling rather than fit to the trained distribution.

Transfer is the discriminating test: we train on one family of forecasting problems and evaluate on
held-out problems the model never saw, reusing the very diagnostics of Experiments~1 and~2 so that an
improvement is legible in the same terms. But transfer alone is not enough, and the design below is
shaped by a single methodological worry. Training against a forecasting reward can lower held-out
error for reasons that have nothing to do with world models: a model can learn the output format, learn
to track the poll's level, or simply learn to stop over-reacting to news. Each of those improves the
score while inferring nothing. Most of Experiment~3's machinery exists to separate that from genuine
latent inference---an instrument that can see the latent, a reward that a shrinkage policy cannot
collect, a bracket of blind references, and a structureless control arm.

\paragraph{Why these two arms.}  The two earlier experiments isolated two different things a world model
must do, and Experiment~3 stresses each under training. Experiment~2 isolated \emph{latent recovery} in
a world we built, where the truth is known by construction; the synthetic arm trains that ability and
asks whether it transfers to worlds withheld from training. Experiment~1 isolated \emph{coherent
updating} on real questions, where the truth is contested and arrives only at resolution; the
pastcasting arm trains on resolved questions and asks whether the coherence signature---not just raw
accuracy---survives on held-out ones.

% ------------------------------------------------------------------
\subsection{Synthetic worlds: training latent-structure inference and testing transfer across structure}
\label{sec:exp3-synthetic}
% ------------------------------------------------------------------

\input{fig_exp3_dag_family}

\textbf{A family of latent-structure worlds.}  We generalize the Experiment~2 world into a
\emph{family} of small linear-Gaussian DAGs that share its logic but vary its structure
(Fig.~\ref{fig:exp3-family}).  In each, observed news $x$ drives one or more latent opinions $o$ that
carry over and mean-revert week to week and emit an observed poll $p$ (Fig.~\ref{fig:exp3-unroll}); a
single edge---the \emph{hidden gain} $g\sim\mathrm{Uniform}(0.1,1.0)$, drawn once per city---is never
revealed.  The $12$ structures vary everything \emph{around} that edge: the number of news drivers, the
depth and delay of the news$\to$poll path, inter-latent couplings and feedback, mediating latents, and
which edge is hidden (an input gain, or the persistence itself).  We split them into $8$ \emph{training}
and $4$ \emph{held-out} structures; among the held-out is the direct chain---the Experiment~2 world
itself---so the test literally asks whether training on other structures improves recovery on the
Experiment~2 world.

Two properties make the split a real test rather than a relabeling. First, the held-out worlds are
\emph{observationally distinct}: applying each structure's Kalman oracle to every other structure's data
(within an interface), every world's own oracle is best by a margin of $21$--$463\%$, so no held-out
world is a disguised copy of a trained one. Second, as in Experiment~2's parity prompt, the generative
structure is \emph{disclosed} in words and only $g$ is hidden: the forecaster is told the graph, the
delays, and the persistences, and must infer from the news/poll history the one latent scalar that
scales news into opinion. Transfer therefore tests whether the model can \emph{use} a
described-but-never-trained-on world model and pin down its free parameter---not whether it memorized a
structure it trained on. It also means the structure-agnostic baselines are \emph{conservative}: the
model is handed strictly more information than they get.

\input{fig_exp3_unroll}
\input{fig_exp3_probe_diagnostic}

\textbf{The instrument must be able to see the latent.}  Experiment~2 reads recovery with a four-shock
predict-then-slope probe: ask for the poll under each of four news shocks and take the slope of the
answers. Carried over naively---asking always about \emph{next} week's poll---that probe is
\emph{degenerate} in any world where news takes more than a week to reach the poll. There, next week's
poll does not depend on next week's news at all: the four scenarios share one target, the true slope is
zero, and a constant reply is exactly optimal. Six of our twelve worlds are blind in precisely this way
(Fig.~\ref{fig:exp3-probe}a). We therefore read each world at its own \textbf{first-response horizon}
$h^\star$: the earliest week at which a news pulse moves the probed poll, i.e.\ the smallest $h$ with
$(C A^{h-1} B)\neq 0$. Under $h^\star$ the probe slope carries $g$ in $6/8$ training and $4/4$ held-out
worlds. The two exceptions are the hidden-$\phi$ worlds, where $g$ is a persistence rather than an input
gain; there no input$\to$output response identifies it at any horizon, correctly and by construction, and
those worlds are scored on forecast error alone.

\textbf{A reward a shrinkage policy cannot collect.}  The obvious reward---accuracy of the predicted
poll---is \emph{improvable without inferring anything}. Because $g$ is drawn around a known mean, a
policy that ignores the news and answers a well-chosen constant already captures most of the available
reward: under the one-step probe, $70\%$ of the oracle's (Fig.~\ref{fig:exp3-probe}b). This is not
hypothetical. An earlier generation of our runs optimized exactly this reward and converged exactly
there: forecast error fell while informativeness \emph{collapsed} ($\rho: .58\!\to\!.26$), the signature
of a model that had learned a better average response rather than a per-city inference. We therefore
score the \emph{response}, not only the level. With the four scenario forecasts $\hat p_i$ elicited
jointly in one reply,
\begin{equation}
r \;=\; (1-w)\cdot\underbrace{\tfrac14\textstyle\sum_i \max\!\big(0,\,1-|\hat p_i-\mu_i|/S\big)}_{\text{level}}
\;+\; w\cdot\underbrace{\max\!\big(0,\,1-|\mathrm{slope}(\hat p)-\gamma^\star|/S_g\big)}_{\text{response}},
\label{eq:exp3-reward}
\end{equation}
where $\mu_i$ is the noise-free target under shock $i$ and $\gamma^\star=(CA^{h^\star-1}B)$ is that
city's true response at $h^\star$ ($w{=}0.5$, $S{=}15$ poll points, $S_g{=}0.5$). The slope term is
per-episode: it can only be earned by reading \emph{this} city's history. It cuts a constant reply from
$70\%$ to $46\%$ of the oracle's reward, and it does so in every world, not only the immediate ones.

\textbf{Three metrics, because one is not enough.}  We read the same three axes as Experiment~2, and we
read them together. $\pi=\overline{|\hat p-\mu|}$ is the next-poll forecast error (poll points; defined
in every world). $\varepsilon=\overline{|\hat\gamma-\gamma^\star|}$ is the recovery error, and
$\rho(\hat\gamma,\gamma^\star)$ the informativeness. Neither $\pi$ nor $\varepsilon$ alone certifies
inference: a good constant response buys low $\pi$ with no inference at all, and shrinking $\hat\gamma$
onto a world's mean response buys low $\varepsilon$ while ordering no cities. Only a forecaster that is
calibrated \emph{and} informative---low $\varepsilon$ with $\rho$ well above zero---is recovering the
latent. We meet a live instance of this below, on the deepest held-out chain.

\textbf{Recovery must be read on each world's own scale.}  The recoverable response
$\gamma^\star=(CA^{h^\star-1}B)$ spans $0.10$--$1.00$ in the direct chain but only $0.022$--$0.216$ in
the four-stage chain, so a raw $\varepsilon$ is not comparable across worlds and must not be pooled.
We therefore normalise by each world's own \emph{blind} floor,
\begin{equation}
\tilde\varepsilon \;=\; \varepsilon \,/\, \varepsilon_{\text{prior-blind}},
\label{eq:exp3-epstilde}
\end{equation}
where prior-blind is that structure's Kalman filter with $g$ frozen at the range midpoint: it knows the
structure and never infers the gain, so its $\varepsilon$ is exactly the cost of \emph{not} inferring,
expressed on the right scale. $\tilde\varepsilon<1$ means inference beat not inferring.

\textbf{References that are actually hard.}  The oracle is a ceiling, not a bar. Because the prompt
discloses the structure, the sharp comparison is against forecasters that also know the structure but
never infer $g$. We compute six references on the \emph{identical} episodes (Table~\ref{tab:exp3-refs}):
the \emph{oracle} (Kalman filter marginalizing $g$); \emph{prior-blind} (the same filter with $g$ frozen
at the range midpoint $0.55$); the \emph{best fixed $\gamma$}, the strongest single constant response
chosen in hindsight per world; \emph{always-prior} ($\gamma{=}0.55$); a structure-agnostic \emph{naive
frequentist}; and \emph{persistence}. The blind constants are the bar a genuine inference must clear.

\begin{table}[h]
\centering
\caption{\textbf{The bracket, per held-out world}, computed on the identical episodes at each world's
$h^\star$ and over its informative prefixes. The oracle is a ceiling, not a bar: the references that
matter are the blind ones. Freezing $g$ at the range midpoint but using the disclosed structure
(prior-blind) already forecasts the direct chain to $2.07$ poll points, and the best single constant
response to $2.36$---so a low $\pi$ is not by itself evidence of inference. Note also how far apart the
worlds are: $\pi$ in the delayed worlds is three to four times that of the immediate ones, and the
recovery floors differ by $4\times$; this is why we neither pool $\pi$ nor pool raw $\varepsilon$.
\emph{Generated by \texttt{make\_tables.py}.}}
\label{tab:exp3-refs}
\footnotesize
\input{tables/exp3_refs}
\end{table}

\textbf{An identifiability guard.}  In a delayed world at a short prefix, the news has not yet reached
the poll and \emph{no} forecaster can identify $g$: the oracle and the prior-blind filter coincide. Cells
like these would dilute any claim, so we score only (structure, prefix) cells where the oracle beats
prior-blind by at least $0.1$ poll points. In practice this keeps every prefix for the immediate worlds
and only $k\geq5$ for the four-stage chain and the mediator graph.

\textbf{The discriminating control.}  Even with all of the above, a held-out improvement could still come
from learning the JSON format, from tracking the poll's level, or from learning to suppress a
news-over-reaction the base model has. None of these is a world model. We therefore train a second arm
on \textbf{structureless} data: the same worlds and the same prompts, but the \emph{displayed} news is
resampled independently of the polls (zero mutual information), every scenario's target is the expected
poll under \emph{zero} input, and the slope target is $0$. A control-trained model can learn format,
level-tracking, and news-suppression---but there is no news$\to$poll structure in its data to learn. Both
arms are then evaluated on the \emph{identical} real held-out worlds. This yields a sharp verdict rule:

\begin{quote}
\emph{Family training shows transfer-as-world-modeling only if, relative to the untrained base, $\pi$
improves, $\pi$ beats the blind constants (always-prior and best-fixed), $\rho$ does not collapse toward
zero---and the family arm separates from the control. A family arm that matches the control
learned format, not structure.}
\end{quote}

\input{fig_exp3_transfer_curves}

\textbf{Setup and status.}  We train Qwen-3-4B (LoRA rank $32$) with Dr.\ GRPO on the reward of
Eq.~\eqref{eq:exp3-reward} over the $8$ training structures, paired seeds per arm, and evaluate on the
$4$ held-out structures with the unchanged Experiment-2 instrument (full configuration in
Appendix~\ref{app:exp3}). The base model is a genuinely poor forecaster on these worlds---it over-reads
the news badly---so there is ample room for a training gain that means nothing. That is precisely what
the control is for. Training is in progress; Table~\ref{tab:exp3-transfer} reports the current
per-world state, Fig.~\ref{fig:exp3-curves} tracks the three axes against the bracket, and
Fig.~\ref{fig:exp3-per-env} resolves them by world.

\textbf{What the current checkpoints show.}  Training improves every world substantially over the
untrained base, and on the two immediate worlds both arms move from well above persistence to well below
it. But the family arm does not separate from the structureless control on the direct chain---the
Experiment-2 world---on either $\pi$ or $\tilde\varepsilon$, and the control is there the \emph{more}
informative of the two. Where the arms do separate is on the delayed worlds, and inconsistently: the
family arm carries real $\rho$ on \texttt{mediators} and \texttt{chain4} where the control's collapses,
yet on \texttt{chain4} the control still posts the lower raw $\varepsilon$ by shrinking onto that
world's mean response. Meanwhile $\tilde\varepsilon\geq0.96$ everywhere: neither arm has yet beaten the
cost of never inferring the gain at all. On the present evidence the large improvement over the
untrained model is attributable to format and level-tracking---what the control also learns---rather
than to recovered structure.

\begin{table}[h]
\centering
\caption{Synthetic arm --- \textbf{transfer across structure, per held-out world}. Qwen-3-4B trained with
Dr.\ GRPO on the $8$ training structures (family) and, identically, on structureless data (control), read
on the $4$ held-out structures with the Experiment~2 four-shock instrument at each world's $h^\star$;
paired seeds, bold marks the better arm on $\pi$. $\tilde\varepsilon$ is the recovery error relative to
that world's blind floor (Eq.~\eqref{eq:exp3-epstilde}): $\tilde\varepsilon<1$ means inferring $g$ beat
\emph{not} inferring it. \textbf{There is deliberately no pooled row}: $\pi$ would be dominated by the
delayed worlds and raw $\varepsilon$ would average incommensurable scales. Two readings matter. First,
$\tilde\varepsilon\geq0.96$ everywhere---neither arm clears the blind floor on any world. Second,
\texttt{chain4} is the reason $\varepsilon$ and $\rho$ must be read together: the control attains the
\emph{lower} raw $\varepsilon$ there while its $\rho$ collapses to $.081$, because it shrinks
$\hat\gamma$ onto that world's small mean response---calibrated but blind, exactly the solution the
reward was built to exclude. \emph{[Runs in flight; generated by \texttt{make\_tables.py}.]}}
\label{tab:exp3-transfer}
\footnotesize
\input{tables/exp3_transfer}
\end{table}

% ------------------------------------------------------------------
\subsection{Real-world pastcasting: training coherent updating and testing it on held-out questions}
\label{sec:exp3-pastcasting}
% ------------------------------------------------------------------

\textbf{Setup.}  The pastcasting arm reuses the Experiment~1 pipeline---the same structured forecast
(event model, competing hypotheses $H_1/H_2$, headline $\hat p = P(H_1)$) and the same agents---but
swaps the \emph{unresolved} prospective markets for \emph{resolved} historical Polymarket questions,
restricting each agent to evidence available at the time. This closes the one gap Experiment~1 had to
defer: with realized outcomes in hand we can compute a genuine accuracy reward (Brier / log-loss against
resolution), the ``resolution follow-up'' flagged in Appendix~\ref{app:exp1-limitations}. We train on
one family of questions and evaluate on a held-out set, so the score measures forecasting skill that
generalizes across questions, not memorization of particular resolutions.

\textbf{Transfer is more than accuracy.}  Experiment~1 established that frontier agents already track
without copying the market and update coherently; the risk in optimizing a resolution reward is that a
model recovers accuracy by reviving exactly the shortcuts Experiment~1 ruled out---market imitation or
surface matching. We therefore evaluate the held-out set on \emph{both} axes. (i) \emph{Accuracy}: Brier
and log-loss on held-out resolved questions, relative to the pre-training agent and to the contemporary
market price. (ii) \emph{Coherence}: we re-run the Experiment~1 counterfactual probe on the held-out
questions and recompute EHC/HFC and the sensitivity ratio
(Section~\ref{sec:exp1-results}). Training should preserve or sharpen the coherent-updating
signature---larger, more causally selective revisions on relevant evidence, near-zero movement on
orthogonal evidence---rather than trade it away for in-distribution fit. Improvement on held-out
\emph{questions}, in the same coherence terms Experiment~1 introduced, is the real-world counterpart of
the held-out-world recovery gain in the synthetic arm.

% ------------------------------------------------------------------
\paragraph{What the two arms together would show.}  The arms are deliberately complementary along the
two dimensions Experiments~1 and~2 separated. The synthetic arm tests whether training improves
\emph{latent-structure recovery} and carries it to an unseen generative process, scored where the truth
is known and against a control that isolates the part of the gain attributable to structure. The real arm
tests whether training improves \emph{coherent updating and accuracy} and carries it to unseen questions,
scored where the truth eventually resolves. Read through the same metrics as the earlier experiments, a
consistent lift on held-out distributions in both arms---one that the control cannot explain---is the
cumulative evidence the diagnostic was built for: agents \emph{have} partial inductive world models,
those models can be \emph{trained}, and the improvement is general.
